% !PRINTED 424
\weberpage{424}

%<<CONT
so folgt nach (1) und \S.\,111, (12):
\weq{7}{Y = \lambda v + \mu vu,}
und wenn wir, wie in \S.\,111
\weq{8}{\frac{H^3}{f^5} = z,\quad \frac{T^2}{f^5} = z_1,\quad z + z_1 = 1728}
setzen, so gehen die Ausdr\"ucke (3), (4), (5) in folgende \"uber:
\weq{9}{az = 8\lambda^3 + \lambda^2\mu + \frac{72\lambda\mu^2 + \mu^3}{z_1}}
\[
bz = -\lambda^4 + \frac{18\lambda^2\mu^2 + \lambda\mu^3}{z_1} + 27\frac{\mu^4}{z_1^2}
\]
\[
cz = \lambda^5 - 10\frac{\lambda^3\mu^2}{z_1} + \frac{45\lambda\mu^4 + \mu^5}{z_1^2},
\]
wenn $Y$ die Wurzel der Gleichung
\weq{10}{Y^5 + 5aY^2 + 5bY + c = 0}
ist. Diese Gleichung hat die Form einer Hauptgleichung 5\textsuperscript{ten} Grades, und sie kann jede Hauptgleichung darstellen, wenn wir $a, b, c$ beliebig annehmen k\"onnen. Man hat also, wenn $a, b, c$ beliebig gegeben angenommen werden, aus (9) die drei Gr\"ossen $\lambda, \mu, z_1$ (und $z = 1728 - z_1$) zu bestimmen, und es ist eine sehr merkw\"urdige Eigenth\"umlichkeit dieses Gleichungssystemes, dass sich diese Unbekannten \gesperrt{rational} durch die gegebenen Gr\"ossen $a, b, c$ und die Discriminante der Gleichung (10) ausdr\"ucken lassen. Um dies nachzuweisen, leiten wir aus (9) zun\"achst einfachere Gleichungen ab.

Wenn wir die zweite von ihnen mit $\lambda$ multipliciren und zur dritten addiren, so folgt
\weq{11}{z_1(\lambda b + c) = \mu^2a,}
und wenn wir die dritte mit $\lambda$, die zweite mit $\mu^2 : z_1$ multipliciren und subtrahiren
\weq{12}{z\left(\lambda c - \frac{\mu^2}{z_1}b\right) = \left(\lambda^2 - 3\frac{\mu^2}{z_1}\right)^3.}

Ebenso ergiebt sich aus der ersten und zweiten
\[
z\,\frac{a\lambda + 8b}{\mu} = \lambda^3 + \frac{216\lambda^2\mu}{z_1} + 9\frac{\lambda\mu^2}{z_1} + \frac{216\mu^3}{z_1^2}.
\]
